% =========================================================================
% Figura: Interpretación geométrica del indicador IGD+
% Problema de minimización bidimensional
% =========================================================================

\begin{figure}[htbp]
\centering
\resizebox{0.58\linewidth}{!}{%
\begin{tikzpicture}[
    font=\sffamily,
    >=Stealth,
    x=1cm, y=1cm
]

% -------------------------------------------------------------------------
% 1. EJES COORDENADOS (Minimización)
% -------------------------------------------------------------------------
\draw[line width=1.1pt, draw=darkText, ->, -{Stealth[length=6pt, width=4pt]}]
    (0, 0) -- (6.8, 0) node[below=5pt, midway, font=\fontsize{9.5}{12}\selectfont, text=darkText] {\textbf{Objetivo $f_1$} (minimizar)};

\draw[line width=1.1pt, draw=darkText, ->, -{Stealth[length=6pt, width=4pt]}]
    (0, 0) -- (0, 5.7);
\node[font=\fontsize{9.5}{12}\selectfont, text=darkText, rotate=90]
    at (-0.65, 2.85) {\textbf{Objetivo $f_2$} (minimizar)};

% -------------------------------------------------------------------------
% 2. FRENTE DE REFERENCIA (Curva suave y puntos de referencia Z en figOrange)
% -------------------------------------------------------------------------
\draw[draw=figOrange!75, dashed, line width=1.0pt]
    (0.9, 4.5) to[out=-70, in=150] (1.8, 2.7) to[out=-30, in=165] (3.8, 1.2) to[out=-15, in=175] (6.0, 0.65)
    node[right=2pt, font=\fontsize{8.5}{10.5}\selectfont\bfseries, text=figOrange!90!black] {$PF_{\text{ref}}$};

% Coordenadas de los puntos de referencia Z (5 puntos bien distribuidos)
\coordinate (z1) at (1.15, 3.8);
\coordinate (z2) at (1.80, 2.7);
\coordinate (z3) at (2.70, 1.8);
\coordinate (z4) at (3.80, 1.2);
\coordinate (z5) at (5.10, 0.8);

% -------------------------------------------------------------------------
% 3. CONJUNTO DE SOLUCIONES OBTENIDAS A (Círculos azules figBlue)
% -------------------------------------------------------------------------
\coordinate (a1) at (1.60, 4.3);
\coordinate (a2) at (2.35, 3.2);
\coordinate (a3) at (3.30, 2.2);
\coordinate (a4) at (4.45, 1.6);
\coordinate (a5) at (5.65, 1.2);

% -------------------------------------------------------------------------
% 4. DISTANCIAS MODIFICADAS d+ (Flechas bermellón figVermilion)
% -------------------------------------------------------------------------
\draw[line width=1.1pt, draw=figVermilion, -{Stealth[length=5pt, width=3.5pt]}] (z1) -- (a1);
\draw[line width=1.1pt, draw=figVermilion, -{Stealth[length=5pt, width=3.5pt]}] (z2) -- (a2);
\draw[line width=1.1pt, draw=figVermilion, -{Stealth[length=5pt, width=3.5pt]}] (z3) -- (a3);
\draw[line width=1.1pt, draw=figVermilion, -{Stealth[length=5pt, width=3.5pt]}] (z4) -- (a4);
\draw[line width=1.1pt, draw=figVermilion, -{Stealth[length=5pt, width=3.5pt]}] (z5) -- (a5);

% Dibujar puntos de referencia Z (círculos naranja figOrange)
\newcommand{\drawZ}[1]{%
    \filldraw[draw=figOrange!90!black, fill=figOrange, line width=0.8pt] (#1) circle (0.10cm);
}
\drawZ{z1}
\drawZ{z2}
\drawZ{z3}
\drawZ{z4}
\drawZ{z5}

% Dibujar soluciones obtenidas A (círculos azules figBlue)
\newcommand{\drawA}[1]{%
    \filldraw[draw=figBlue!90!black, fill=figBlue, line width=0.9pt] (#1) circle (0.12cm);
}
\drawA{a1}
\drawA{a2}
\drawA{a3}
\drawA{a4}
\drawA{a5}

% -------------------------------------------------------------------------
% 5. LEYENDA (Esquina superior derecha libre)
% -------------------------------------------------------------------------
\begin{scope}[shift={(2.90, 3.75)}]
    \draw[draw=grayText!40, fill=white, rounded corners=3pt, line width=0.6pt]
        (0, 0) rectangle (3.75, 1.30);

    % Soluciones obtenidas A
    \filldraw[draw=figBlue!90!black, fill=figBlue, line width=0.6pt] (0.20, 1.00) circle (0.08cm);
    \node[anchor=west, font=\fontsize{6.5}{7.8}\selectfont, text=darkText]
        at (0.36, 1.00) {Soluciones obtenidas ($A$)};

    % Puntos de referencia Z
    \filldraw[draw=figOrange!90!black, fill=figOrange, line width=0.6pt] (0.20, 0.65) circle (0.07cm);
    \node[anchor=west, font=\fontsize{6.5}{7.8}\selectfont, text=darkText]
        at (0.36, 0.65) {Puntos de referencia ($Z$)};

    % Distancia modificada d+
    \draw[line width=1.0pt, draw=figVermilion, -{Stealth[length=4pt, width=2.8pt]}]
        (0.08, 0.30) -- (0.32, 0.30);
    \node[anchor=west, font=\fontsize{6.5}{7.8}\selectfont, text=darkText]
        at (0.36, 0.30) {Distancia modificada ($d^+$)};
\end{scope}

\end{tikzpicture}%
}
\caption{Interpretación geométrica del indicador de distancia generacional invertida modificada (\acrshort{igd}$^+$) en un espacio de dos objetivos a minimizar. Desde cada punto del frente de referencia ($Z$) se calcula la distancia más cercana ($d^+$) al conjunto de soluciones obtenidas ($A$).}
\label{fig:indicador_igd_plus}
\end{figure}
